\chapter{Potential gradient descent and baryonification}\label{ch:pgd-baryonification}
\chapterstatus{written}{2026-11-20}

Two post-processing methods move the particles of a gravity-only simulation to mimic physics it does not resolve or
contain: potential gradient descent (PGD) sharpens the haloes of low-resolution PM simulations, and baryonification
displaces particles around haloes to mimic the effect of gas, stars and feedback on the matter distribution.

\begin{modulebox}{PGD and Baryonification}
\file{PGD/PotentialGradientDecent.h} (example \file{PGD/example}); \file{Baryonification/Baryonification.h},
\file{Baryonification/Profiles.h} (namespace \nsp{FML::BARYONIFICATION}); SimulationCode integration
\file{src/BaryonificationOutput.h/.cpp}.\par
Documentation: \file{Baryonification/README.md}, \file{SimulationCode/ANALYSIS.md}.
\end{modulebox}

\section{Potential gradient descent}\label{sec:pgd}

PM simulations with few steps and coarse force grids produce haloes that are too extended. Dai, Feng and Seljak
\cite{Dai2018PGD} move the particles a little further down the gradient of a band-pass filtered potential:
\begin{equation}
\bm S=-\alpha\,\mathcal N\,\nabla\,\hat O_l\hat O_s\nabla^{-2}\delta,\qquad
\hat O_l=e^{-k_l^2/k^2},\qquad \hat O_s=e^{-k^4/k_s^4},
\end{equation}
where $\hat O_l$ removes the large scales (which the PM simulation gets right) and $\hat O_s$ the scales below the
resolution. With the paper's definition of the fitted, dimensionless $\alpha(a)$ the prefactor is
$\mathcal N=4\pi G\bar\rho/H_0^2=\frac32\Om_m$ (with $\bar\rho$ the comoving mean density); a positive $\alpha$ moves
particles towards overdensities. The density is assigned to a grid, filtered and differentiated in Fourier space, and
the displacement interpolated back to the particles. The function is collective; the \cls{MPIParticles} overload
communicates the moved particles. Parameters: $k_l$, $k_s$ (in box units) and $\alpha$, calibrated by the paper's
fitting formulae (\code{pgd_fitting_formula_*}).

\section{Baryonification}\label{sec:baryonification}

\paragraph*{Model.} Following Schneider and Teyssier \cite{SchneiderTeyssier2015} and the profile family of
Schneider et al.\ \cite{Schneider2019}, each halo's gravity-only (reference) profile is replaced by a profile of
collisionless matter, gas and a central galaxy, and particles are displaced so that enclosed masses match. With
$x=r/R_{200}$:
\begin{itemize}
\item reference: a truncated NFW profile $\propto[cx(1+cx)^2(1+(x/\epsilon)^2)^2]^{-1}$, normalised to $M_{200}$;
\item gas: $\propto(1+x/\theta_{\rm co})^{-\beta}[1+(x/\theta_{\rm ej})^2]^{-(7-\beta)/2}$ with $\beta=3-(M_c/M_{200})^\mu$;
\item central galaxy: $\rho\propto r^{-2}\exp[-(r/2R_h)^2]$ (cumulative mass analytic with erf);
\item collisionless matter (dark matter and satellite stars), optionally relaxed by
      $r_f/r_i-1=A[(M_i/M_f)^n-1]$, $M_f=f_{\rm clm}M_i+M_{\rm gas}(r_f)+M_{\rm cga}(r_f)$ (bisection);
\item stellar fractions $f_\star=A_\star(M_1/M_{200})^{\eta}$, gas $f_b-f_\star$.
\end{itemize}
\paragraph*{Displacement.} With the cumulative masses $M_{\rm ref}(r)$ and $M_{\rm bar}(r)$, a particle at $r$ moves
radially to $r'$ with $M_{\rm bar}(r')=M_{\rm ref}(r)$ (\cls{RadialMapping} of monotone piecewise-linear tables). All
components are normalised at a finite support (default $30R_{200}$) with a uniform background, so the displacement
vanishes outside it; overlapping haloes add their displacement vectors, evaluated at the original positions. This is a
finite-support implementation of the profile family, not a reproduction of the paper's calibrated predictions; there is
no correlated two-halo term. Velocities, masses and IDs are unchanged.

\paragraph*{In SimulationCode.} \code{baryonification = true} runs the method at each output on a copy of the particles
(the live simulation is unchanged): FoF seeds, spherical-overdensity $M_{200c}$ and $R_{200c}$ with shrinking-sphere
centres, rejection of truncated, nested or unresolved hosts, a power-law concentration (illustrative, not fitted from PM
cores), and outputs of the corrected and original $P(k)$ and their ratio, the host catalogue and optionally the
displaced particles. Requires 3D GR, one CDM+baryon species and no massive neutrinos.

The parameters map onto the model above (\param{baryonification_model} \code{= schneider2019}, the only family; the
defaults are illustrative, not calibrated). Gas: $M_c$ (\param{baryonification_gas_mass_scale_msunh}), $\mu$
(\param{baryonification_gas_mass_exponent}), $\theta_{\rm co}$ and $\theta_{\rm ej}$ in units of $R_{200c}$
(\param{baryonification_gas_core_radius_over_r200}, \param{baryonification_gas_ejection_radius_over_r200}); the NFW
truncation $\epsilon$ (\param{baryonification_truncation_radius_over_r200}). Stars: $M_1$, $A_\star$, the exponents of the
total and central stellar fractions and $R_h/R_{200c}$ (\param{baryonification_stellar_mass_scale_msunh},
\param{baryonification_stellar_amplitude}, \param{baryonification_stellar_total_exponent},
\param{baryonification_stellar_central_exponent}, \param{baryonification_stellar_radius_over_r200}). Relaxation:
\param{baryonification_relax_collisionless} with $A$ and $n$ (\param{baryonification_relaxation_amplitude},
\param{baryonification_relaxation_exponent}). Concentration $c=c_0(M_{200c}/M_{\rm piv})^{\alpha}(1+z)^{\gamma}$
(\param{baryonification_concentration_norm}, \param{baryonification_concentration_pivot_mass_msunh},
\param{baryonification_concentration_mass_slope}, \param{baryonification_concentration_redshift_power}). Hosts: FoF seeds
with \param{baryonification_fof_linking_length}, \param{baryonification_fof_min_particles},
\param{baryonification_fof_nmesh_max} and \param{baryonification_fof_buffer_mpch} (larger than the halo radii), and
$M_{200c}\ge$\,\param{baryonification_min_halo_mass_msunh}. Tables: \param{baryonification_profile_radial_bins} radii from
\param{baryonification_min_radius_over_r200} to the support \param{baryonification_max_radius_over_r200} (less than half the
box). \param{baryonification_output_particles} writes the displaced particles as a separate snapshot.

\begin{cppcode}[Baryonifying a catalogue]
#include <FML/Baryonification/Baryonification.h>
using namespace FML::BARYONIFICATION;
HaloProperties halo{1e14, 0.7, 5.0};               // M200 [Msun/h], R200 [Mpc/h], concentration
ModelParameters model;
model.gas_mass_scale = 1e14;  model.gas_ejection_radius = 4.0;
auto profiles = make_halo_profiles(halo, model, /*mean_density=*/8.3e10);
std::vector<HaloDisplacement> halos{{{50.0, 50.0, 50.0}, profiles.mapping()}};
Baryonifier correction(std::move(halos), /*boxsize=*/100.0);
correction.apply(particles.data(), particles.size());   // positions only
\end{cppcode}

\begin{caveat}
Validate against a chosen hydrodynamical calibration, halo completeness, halo mass and concentration errors, snapshot
resolution, the support and the radial tables; no percent-level accuracy is claimed. The halo catalogue and prepared
mappings are replicated on every rank. With PM resolution, SO masses are biased low unless $R_{200}$ spans several force
cells (\cref{ch:simulation-code}).
\end{caveat}
